\section{Gravity, clocks and station observations from one state}
\label{r14-gravity-time}

The satellite specialization shares the orbit, gravitational coefficients,
station motion, frame realization and clock parameters with its observation
map. This closes a specific gap in R11: its orbital station query is
simultaneous geometry, whereas its live GPS query uses emission and reception
events. The following equations define their common physical interface.

\subsection{A source potential with explicit finite exterior models}
For an isolated Newtonian mass density $\rho_m\geq0$, fix the potential zero
at infinity. In SI units the source law and its boundary realization are
\begin{equation}
 \nabla^2\Phi=4\pi G\rho_m,\qquad
 \Phi(x)=-G\int\frac{\rho_m(x')}{\|x-x'\|}\,d^3x',\qquad
 g=-\nabla\Phi .
 \label{r14-poisson}
\end{equation}
Outside the source $\nabla^2\Phi=0$. Boundary conditions or measured exterior
coefficients are required; the Poisson equation alone does not identify them.
\href{https://ocw.mit.edu/courses/8-902-astrophysics-ii-fall-2004/a640fdbf6840889b97a7c3680cb2e4f1_lec4.pdf}{[GT1: MIT, potential theory]}

Let $R(t)\in SO(3)$ map Earth-fixed components to a common geocentric,
nonrotating spatial chart. For $p=R^Tr$, $s=\|p\|>a_E>0$ and $\mu_E>0$,
one closed exterior choice is
\begin{equation}
 \Phi_E(t,r)=-\frac{\mu_E}{s}
 \left[1-J_2\left(\frac{a_E}{s}\right)^2P_2(p_z/s)\right],
 \qquad P_2(u)=\frac{3u^2-1}{2} .
 \label{r14-exterior-potential}
\end{equation}
$J_2=0$ selects the monopole. Alternatively select the finite R11 solid-harmonic
body with $\Phi_E=-U_{\rm R2}$, retaining its normalization, degree, coefficient
epoch and exterior domain. These are alternative Earth fields; adding them
would count the monopole twice. $\mu_E=GM_E$ is a gravitational parameter with
a declared realization, not an exact SI defining constant. Arbitrary fitted
coefficients need not correspond to a nonnegative physical density; consistency
with the selected source and coefficient uncertainty is a separate obligation.

For a point body at geocentric $r_b(t)\ne0$, the appropriate geocentric
\emph{tidal} potential and its force are
\begin{align}
 \Phi_b^{\rm tid}(t,r)
 &=-\frac{\mu_b}{\|r_b-r\|}+\frac{\mu_b}{\|r_b\|}
       +\frac{\mu_b\,r_b\cdot r}{\|r_b\|^3},\\
 -\nabla_r\Phi_b^{\rm tid}
 &=\mu_b\left(\frac{r_b-r}{\|r_b-r\|^3}
                   -\frac{r_b}{\|r_b\|^3}\right),\qquad r\ne r_b .
 \label{r14-tidal-potential}
\end{align}
The constant and linear terms remove the external potential and acceleration
at the geocenter. Set $\Phi_{\rm geo}=\Phi_E+\sum_b\Phi_b^{\rm tid}$ and
\begin{equation}
 \dot r=v,\qquad \dot v=-\nabla_r\Phi_{\rm geo}+a_{\rm ng}+a_{\rm 1PN} .
 \label{r14-orbit-closure}
\end{equation}
$a_{\rm ng}$ binds the selected radiation-pressure, thrust or other physical
force body and its parameters. Set absent terms to zero explicitly. The
existing R11 spherical-Earth $a_{\rm GR}$ may bind $a_{\rm 1PN}$; it is a
specified post-Newtonian correction, not the whole relativistic force model.
All retained force choices remain subject to their source/model error.

\subsection{Frame acceleration and gravitational energy}
For a general Newtonian inertial chart $x=o(t)+R(t)p$, define
$\Omega=R^T\dot R$ and $\Omega z=\omega\times z$. Differentiating twice gives
\begin{align}
 \dot x&=\dot o+R(\dot p+\omega\times p),\\
 \ddot p&=R^T(\ddot x-\ddot o)-2\omega\times\dot p
           -\dot\omega\times p-\omega\times(\omega\times p).
 \label{r14-frame-acceleration}
\end{align}
Thus a rotating-frame propagation includes translation, Coriolis, Euler and
centrifugal terms. In \eqref{r14-orbit-closure} the geocenter acceleration has
already been removed by the tidal potential; it is not subtracted a second
time. The identity follows from $R^TR=I$, which also makes $\Omega$ skew.
Independently fitted matrix entries in R11 do not establish this identity;
an R14 frame realization must prove it or retain a bounded frame discrepancy.
\href{https://ocw.mit.edu/courses/16-07-dynamics-fall-2009/resources/mit16_07f09_lec08/}{[GT2: MIT, rotating axes]}

For the Newtonian part and constant satellite mass $m$, differentiation yields
\begin{equation}
 \frac{d}{dt}\left(\frac{m\|v\|^2}{2}+m\Phi_{\rm geo}(t,r)\right)
 =m\,v\cdot a_{\rm ng}+m\,\partial_t\Phi_{\rm geo} .
 \label{r14-orbital-energy}
\end{equation}
This equation assumes $a_{\rm 1PN}=0$; otherwise its work must be retained and
a relativistic energy model selected for a relativistic conservation claim.
Moving external sources exchange energy through $\partial_t\Phi_{\rm geo}$.
It is not heat or arithmetic drift. An Earth-fixed constant field can still
be time dependent in the nonrotating chart. A finite integration operator
needs its own energy identity; the continuous balance does not certify RK4.

\subsection{Proper time and the same gravitational state}
Take $t$ as TCG and the spatial/force data as compatible geocentric coordinates.
In the weak-field, slow-motion regime the leading proper-time rate of an
ideal clock on trajectory $r_A(t)$ is
\begin{equation}
 \frac{d\tau_A}{dt}=1+
 \frac{\Phi_{\rm geo}(t,r_A)-\|v_A\|^2/2}{c^2}
 +\epsilon_{\rm clk,A}(t) .
 \label{r14-proper-clock}
\end{equation}
$\epsilon_{\rm clk,A}$ accounts for omitted higher-order/model terms; it is
not an arbitrary control input. The negative-potential convention fixes the
sign: a stationary clock deeper in a monopole field has a smaller rate.
The tidal potential above, rather than the full absolute Sun/Moon potential,
is appropriate to this geocentric convention. A barycentric clock model uses
its own coordinate transformation and potential.
\href{https://iers-conventions.obspm.fr/content/chapter10/tn36_c10.pdf}{[GT3: IERS, chapter 10, equations 10.6--10.9]}

For a scaled coordinate $T=\alpha t+\beta$, $\alpha>0$, the chain rule gives
$d\tau_A/dT=(d\tau_A/dt)/\alpha$. TT uses $\alpha=1-L_G$, with exact
$L_G=6.969290134\times10^{-10}$; station, orbit and
gravitational parameter values must be compatible with the chosen coordinate
scale. GPST tags require their declared realization and mapping. A stored
rational timestamp does not prove a physical clock exact or turn
\eqref{r14-proper-clock} into R11's rational time-conversion template.

The readout is modeled by $C_A(t)=t+\delta_A(t)$, where a declared oscillator
and steering model connects $C_A$ to $\tau_A$. Clock initial offset, intrinsic
frequency error and steering belong to the shared uncertainty state. A
broadcast clock product may already contain physical corrections: apply its
documented convention once, rather than adding \eqref{r14-proper-clock} again.
\href{https://gssc.esa.int/navipedia/index.php/Relativistic_Clock_Correction}{[GT4: ESA, relativistic clock correction]}

\subsection{Causal reception, range and pointing}
For satellite emission at $t_e$ and station reception at $t_r>t_e$, define
\begin{equation}
 d=r_s(t_e)-r_g(t_r),\quad \rho=\|d\|>0,\qquad
 t_r-t_e=\frac{\rho}{c}+\Delta_{\rm grav}+\Delta_{\rm med} .
 \label{r14-light-time}
\end{equation}
This is a two-event equation in one coordinate chart. For the leading static
Earth monopole propagation correction, let $s_e=\|r_s(t_e)\|$,
$s_r=\|r_g(t_r)\|$ and use
\begin{equation}
 \Delta_{\rm grav}=\frac{2\mu_E}{c^3}
 \log\frac{s_e+s_r+\rho}{s_e+s_r-\rho} .
 \label{r14-shapiro}
\end{equation}
It requires $s_e+s_r>\rho$ and a nonocculted exterior ray in the weak-field
domain. $\Delta_{\rm med}$ binds a frequency-dependent propagation model
appropriate to the measured signal, or is zero in the declared vacuum model.
Higher gravity terms, bending, atmospheric error and multipath remain in the
observation discrepancy where omitted.
\href{https://iers-conventions.obspm.fr/content/chapter11/tn36_c11.pdf}{[GT5: IERS, chapter 11, ranging equation 11.17]}

For the flat vacuum choice $\Delta_{\rm grav}=\Delta_{\rm med}=0$ and fixed
$t_r$, put $F(u)=u+\|r_s(u)-r_g(t_r)\|/c-t_r$. On a collision-free interval
with $\|\dot r_s\|<c$, $F'(u)=1+\hat d\cdot\dot r_s/c>0$.
A bracket $F(a)\leq0\leq F(b)$ therefore specifies a unique emission time.
Additional delays require a corresponding monotonicity bound; a fixed number
of correction iterations is not identical by definition to this root.

When positions are supplied in Earth-fixed components and both charts share
the geocenter, the retarded geometric vector in reception axes is
\begin{equation}
 d_E=R(t_r)^TR(t_e)r_{s,E}(t_e)-r_{g,E}(t_r),\qquad
 z=L_g(t_r)d_E=(E,N,U)^T .
 \label{r14-retarded-station}
\end{equation}
$L_g\in SO(3)$ is the declared local ENU rotation. Then
$\mathrm{az}=\operatorname{atan2}(E,N)$ and
$\mathrm{el}=\operatorname{atan2}(U,\sqrt{E^2+N^2})$ on their nonsingular
domains. This product includes Earth rotation during flight; adding a second
Sagnac correction would duplicate it. These are retarded geometric angles.
Physical detector pointing additionally binds aberration, ray bending,
refraction and mounting terms when required.
\href{https://gssc.esa.int/navipedia/index.php/Satellite_Coordinates_Computation}{[GT6: ESA, reception-frame coordinates]}

For code readouts $C_r,C_s$, a corresponding pseudorange observation is
\begin{equation}
 P=\rho+c(\Delta_{\rm grav}+\Delta_{\rm med})
       +c[\delta_r(t_r)-\delta_s(t_e)]+d_P+\eta_P .
 \label{r14-pseudorange}
\end{equation}
Here $d_P$ is the remaining modeled discrepancy/bias and $\eta_P$ the declared
measurement error. Orbit coefficients, both clock states, station motion,
Earth orientation and propagation parameters enter the joint uncertainty
set of Section~\ref{r13-measurement-core}. Their correlations are preserved.
Changing gravity, time or frame in one branch changes every dependent branch.
Literal one-bit encoding preserves these equations and their operands; a
physical tolerance still requires the bound in
Equation~\eqref{r13-measurement-total-bound} for the chosen output and interval.
